% Part: methods
% Chapter: proofs
% Section: example-2

\documentclass[../../../include/open-logic-section]{subfiles}

\begin{document}

\olfileid{mth}{prf}{ex2}

\olsection{مثالی دیگر}

\begin{prop}
اگر $A \subseteq C$، آنگاه $A \cup (C \setminus A) = C$.
\end{prop}

\begin{proof}
  فرض کنید $A \subseteq C$. می‌خواهیم نشان دهیم
  $A \cup (C \setminus A) = C$.
  \begin{quote}
    نخست توجه می‌کنیم که این یک گزارهٔ شرطی است. سور کلی آن به‌طور
    ضمنی در نظر گرفته شده است: حکم برای همهٔ مجموعه‌های $A$ و $C$
    برقرار است. پس $A$ و $C$ متغیرهایی برای مجموعه‌های دلخواه‌اند.
    برای اثبات چنین گزاره‌ای، مقدم را فرض و تالی را اثبات می‌کنیم.

    کار را با استفاده از فرض $A \subseteq C$ ادامه می‌دهیم. تعریف
    ~$\subseteq$ را باز کنیم: این فرض یعنی همهٔ !!{element}های~$A$،
    !!{element}های~$C$ نیز هستند. این نکته را می‌نویسیم، زیرا در سراسر
    اثبات به کارمان می‌آید.
  \end{quote}
  بنا بر تعریف~$\subseteq$، چون $A \subseteq C$، برای هر $z$، اگر
  $z \in A$، آنگاه $z \in C$.
  \begin{quote}
    همهٔ تعریف‌های موجود در فرض را باز کرده‌ایم. اکنون می‌توانیم به
    نتیجه بپردازیم. می‌خواهیم نشان دهیم $A \cup (C \setminus A) = C$؛
    پس اثبات را مانند مثال پیش سامان می‌دهیم: نشان می‌دهیم هر
    !!{element} از $A \cup (C \setminus A)$، !!a{element} از~$C$ نیز
    هست و، برعکس، هر !!{element} از $C$، !!a{element} از
    $A \cup (C \setminus A)$ است. می‌توانیم این را کوتاه‌تر چنین
    بنویسیم: $A \cup (C \setminus A) \subseteq C$ و
    $C \subseteq A \cup (C \setminus A)$. اینجا خلاف جهت بازکردن یک
    تعریف حرکت می‌کنیم، اما خواندن اثبات کمی آسان‌تر می‌شود. چون این
    یک عطف است، باید هر دو بخش را اثبات کنیم. برای نشان‌دادن بخش نخست،
    یعنی اینکه هر !!{element} از $A \cup (C \setminus A)$،
    !!a{element} از~$C$ نیز هست، برای $z$ دلخواه فرض می‌کنیم
    $z \in A \cup (C \setminus A)$ و نشان می‌دهیم $z \in C$. از تعریف
    $\cup$ می‌توانیم از $z \in A \cup (C \setminus A)$ نتیجه بگیریم
    که یا $z \in A$ یا $z \in C \setminus A$. تا اینجا باید دیگر با
    این شیوه آشنا شده باشید.
  \end{quote}
  $A \cup (C \setminus A) = C$ اگر و تنها اگر
  $A \cup (C \setminus A) \subseteq C$ و
  $C \subseteq A \cup (C \setminus A)$. نخست اثبات می‌کنیم
  $A \cup (C \setminus A) \subseteq C$. فرض کنید
  $z \in A \cup (C \setminus A)$. پس یا $z \in A$ یا
  $z \in (C \setminus A)$.
  \begin{quote}
    به یک فصل رسیده‌ایم و می‌خواهیم از آن $z \in C$ را اثبات کنیم.
    این کار را با تفکیک حالت‌ها انجام می‌دهیم.
  \end{quote}
  حالت ۱: $z \in A$. چون برای هر $z$، اگر $z \in A$ آنگاه
  $z \in C$، نتیجه می‌شود $z \in C$.
  \begin{quote}
    اینجا از واقعیتی استفاده کردیم که پیش‌تر از فرض حکم، یعنی
    $A \subseteq C$، به دست آمد. حالت نخست کامل است و به حالت دوم،
    $z \in (C \setminus A)$، می‌پردازیم. به یاد آورید که
    $C \setminus A$ نشان‌دهندهٔ \emph{تفاضل} دو مجموعه است؛ یعنی
    مجموعهٔ همهٔ !!{element}های~$C$ که !!{element}های~$A$ نیستند. اما
    هر !!{element} از $C$ که در~$A$ نباشد، به‌ویژه !!a{element} از~$C$
    است.
  \end{quote}
  حالت ۲: $z \in (C \setminus A)$. این یعنی $z \in C$ و
  $z \notin A$. پس به‌ویژه $z \in C$.
  \begin{quote}
    بسیار خوب، جهت نخست را اثبات کردیم. اکنون به جهت دوم می‌پردازیم.
    اینجا $C \subseteq A \cup (C \setminus A)$ را اثبات می‌کنیم. پس
    فرض می‌کنیم $z \in C$ و نشان می‌دهیم
    $z \in A \cup (C \setminus A)$.
  \end{quote}
  اکنون فرض کنید $z \in C$. می‌خواهیم نشان دهیم یا $z \in A$ یا
  $z \in C \setminus A$.
  \begin{quote}
    چون همهٔ !!{element}های $A$، !!{element}های $C$ نیز هستند و
    $C \setminus A$ مجموعهٔ همهٔ چیزهایی است که !!{element}های $C$
    هستند اما در $A$ نیستند، نتیجه می‌شود $z$ یا در $A$ است یا در
    $C \setminus A$. اگر از پیش ندانید چرا نتیجه درست است، شاید این
    بیان کمی مبهم باشد. بهتر است آن را گام‌به‌گام اثبات کنیم. واقعیت
    ساده‌ای که می‌توانیم بدون اثبات به کار ببریم سودمند است:
    $z \in A$ یا $z \notin A$. این را «اصل طرد شق ثالث» می‌نامند: برای
    هر گزارهٔ~$p$، یا $p$ صادق است یا نقیض آن صادق است. اینجا $p$
    همان گزارهٔ $z \in A$ است. چون این یک فصل است، می‌توانیم باز هم
    اثبات به تفکیک حالت‌ها را به کار ببریم.
  \end{quote}
  یا $z \in A$ یا $z \notin A$. در حالت نخست،
  $z \in A \cup (C \setminus A)$. در حالت دوم، $z \in C$ و
  $z \notin A$؛ پس $z \in C \setminus A$. بنابراین در این حالت نیز
  $z \in A \cup (C \setminus A)$.
  \begin{quote}
    اثبات کامل است: نشان دادیم $A \cup (C \setminus A) = C$.
  \end{quote}
\end{proof}

\end{document}
